\shorthandoff{<>}% Babel español trata como abreviaturas los delimitadores de curvatura de Xy-pic.

\Subsection*{1) Caso de los haces de conjuntos constantes}

1) 1. Reducción al caso en que $X$ es normal sobre~$S$. Sea
$X_{1\overline{s}}$ la normalización de~$(X_{\overline{s}})_{\red}$;
tras restringir $S$ a un abierto no vacío y efectuar una extensión
radicial de~$S$, puede suponerse que $X_{1\overline{s}}$ procede de un
esquema $X_1$ normal sobre~$S$, y que el morfismo
$X_{1\overline{s}}\to X_{\overline{s}}$ procede de un morfismo finito
sobreyectivo $p\colon X_1\to X$ (EGA~IV 8.8.2 y 9.6.1). Supongamos
demostrado el teorema para $fp$. Tras restringir $S$ a un abierto,
puede suponerse que, para todo haz de conjuntos constante finito $F$
sobre~$X$, $(p^*F,fp)$ es cohomológicamente propio relativamente a~$S$
en dimensión $\leqslant 0$ y que $f_*p_*(p^*F)$ es constructible. Según
\Ref{XIII.1.9}, $(p_*p^*F,f)$ es entonces cohomológicamente propio
relativamente a~$S$ en dimensión $\leqslant 0$. El morfismo
$$
F \to p_* p^* F=G
$$
es un monomorfismo. De ello resulta ya, pues $f_*F$ es un subhaz
de~$f_*G$, que $f_*F$ es constructible (SGA~4 IX~2.9~(ii)) y que
$(F,f)$ es cohomológicamente propio relativamente a~$S$ en dimensión
$\leqslant -1$.

Sean $X_2=X_1 \times_{X} X_1$ y $q \colon X_2 \to X$ el morfismo
canónico. Según \Ref{XIII.1.11}~1),
\marginpar{397}
se tiene una sucesión exacta
$$
\xymatrix@C=.5cm{F \ar[r] & G \ar@<2pt>[r]\ar@<-2pt>[r] & q_*q^*F }.
$$

Por lo que acaba de demostrarse, aplicado a $fq$ en lugar de~$f$, puede
suponerse, tras restringir $S$ a un abierto no vacío, que, para todo haz
de conjuntos constante finito $F$ sobre~$X$, $(q^*F,fq)$ es
cohomológicamente propio relativamente a~$S$ en dimensión $\leqslant
-1$, y por tanto que $(q_*q^*F,f)$ es cohomológicamente propio
relativamente a~$S$ en dimensión $\leqslant -1$. Resulta entonces de
\Ref{XIII.1.13}~1) que $(F,f)$ es cohomológicamente propio relativamente
a~$S$ en dimensión $\leqslant 0$.

1) 2. Reducción al caso en que $X$ es normal y afín sobre~$S$.

Sea $U_s$ un abierto afín de~$X_s$, denso en $X_s$. Tras restringir $S$
a un abierto no vacío, puede suponerse que $U_s\to X_s$ se levanta a una
inmersión abierta $i\colon U\to X$, esquemáticamente dominante
relativamente a~$S$ (EGA~IV 8.9.1). Como el morfismo $X\to S$ es
normal, según SGA~2 XIV~1.18 se tiene:
$$
\ifthenelse{\boolean{orig}}
{\prof \et_{S- U} (X) \ge 2 \quoi ;}
{\prof \et_{X- U} (X) \ge 2 \quoi ;}
$$
por consiguiente, para todo haz de conjuntos constante finito $F$
sobre~$X$, el morfismo canónico
$$
F \to i_*i^*F
$$
es un isomorfismo. De ello resulta que, si se supone demostrado el
teorema para $fi$ y para~$i$, entonces, tras restringir $S$ a un abierto
no vacío, $(i^*F,i)$ y $(i^*F,fi)$ son cohomológicamente propios
relativamente a~$S$ en dimensión $\leqslant 0$. Lo mismo vale, pues,
para~$(F,f)$ (\Ref{XIII.1.6}~2)). Como, además,
$f_*F=(fi)_*(i^*F)$ es constructible, queda terminada la reducción.

1) 3. Fin de la demostración.

Puede suponerse que $S$ es normal (EGA~IV 7.8.3). Puede encontrarse una
compactificación de~$X_s$:
\marginpar{398}
$$
\xymatrix{ X_s \ar[r]^{j_s} \ar[d]_{f_s} & P_s \ar[ld]^{g_s} \\ Y_s }
$$
donde $j_s$ es una inmersión abierta dominante y $g_s$ un morfismo
propio; tras efectuar una extensión radicial de~$\kres(s)$ y sustituir
$P_s$ por su normalización, lo que no modifica $X_s$, puede suponerse
que $P_s$ es geométricamente normal. Tras restringir $S$ a un abierto
no vacío y efectuar una extensión radicial sobreyectiva, puede suponerse
que el diagrama anterior procede de un diagrama
$$
\xymatrix{ X \ar[r]^{j} \ar[d]_{f} & P \ar[ld]^{g} \\ Y }
$$
donde $P$ es un esquema normal sobre~$S$, $j$ una inmersión abierta
esquemáticamente dominante relativamente a~$S$ y $g$ un morfismo propio
(EGA~IV 6.9.1, 9.9.4 y 9.6.1). Para todo haz de conjuntos constante
finito $F$ sobre~$X$, de valor $C$, $j_*F$ es el haz constante de
valor~$C$ (SGA~4 \ifthenelse{\boolean{orig}}{}{IX~}2.14.1), y lo mismo
ocurre tras todo cambio de base $S'\to S$;
de ello resulta que $(F,j)$ es cohomológicamente propio relativamente
a~$S$ en dimensión $\leqslant 0$. Lo mismo vale, pues, para~$(F,f)$,
puesto que $g$ es propio \eqref{XIII.1.8}. Como $g$ es propio,
$f_*F=g_*C_P$ es constructible, lo que termina la demostración de
1)\kern2pt a).
